\section{从回答问题到参与科学发现}

这堂课讨论的不是怎样让一个聊天模型写得更像论文，而是怎样把 foundation model 组织成一套\textbf{可与科学家共同工作的发现系统}。主讲人 Vivek Natarajan 把目标概括为加快科学发现的“时钟”：系统不只检索已有知识，还要生成可检验假设、比较彼此冲突的解释、提出实验路径，并把最值得占用实验资源的候选交给人类专家。

\subsection{“加速时钟”究竟在加速什么}

科学工作的耗时并不只来自做实验。研究者还要完成文献筛选、跨学科联想、因果机制构造、候选排序、失败解释和下一轮实验设计。若 AI 只把每一步写得更快，却不能提高\textbf{被验证的新知识}产出，时钟并没有真正加速。因此，本讲把速度与证据绑定，而不是把 token throughput 当作科研产能。

\lecturefigure{slide-002.jpg}{AI co-scientist 的任务：加速科学发现的时钟}{Stanford CS25 V6 Lecture 07 官方录像 00:01:46}

\begin{knowledgebox}{读图：这里的 clock speed 不是模型延迟}
绿色强调的是 scientific discovery，而不是每秒生成多少字。真正应比较的是：从研究问题出现，到形成候选假设、完成关键实验、排除错误解释、得到可复核结论，整个闭环需要多久。LLM 可以压缩阅读和写作时间，但只有当它帮助产生更好的实验决策并缩短验证周期时，才改变科学时钟。
\end{knowledgebox}

可以用一个粗略的发现速率来约束讨论：
$$
R_{\text{discovery}}
=\frac{\Delta K_{\text{validated}}}
{T_{\text{human}}+T_{\text{compute}}+T_{\text{experiment}}}.
$$
其中：
\begin{itemize}
  \item $\Delta K_{\text{validated}}$ 表示新增且被外部证据支持的知识，而不是生成的假设数量；
  \item $T_{\text{human}}$ 是专家阅读、判断、协调与解释时间；
  \item $T_{\text{compute}}$ 是搜索、推理、工具调用和 agent 协作时间；
  \item $T_{\text{experiment}}$ 是湿实验、复现实验或真实世界验证时间。
\end{itemize}

\begin{importantbox}{AI co-scientist 的工作定义}
AI co-scientist 是一个 scientist-in-the-loop 的多智能体系统：科学家给出研究目标、约束、偏好与已有想法，系统生成、批判、比较和演化假设，并把证据、实验建议与不确定性组织成可供人类决策的研究提案。它的交付物是\textbf{待验证的研究行动候选}，不是自动成立的科学真理。
\end{importantbox}

\teachervoice{00:01:06--00:02:20，老师强调目标是给科学家“superpowers”，让 AI 成为 collaborative partner。这里的主语仍是科学家；系统价值来自增强研究者的搜索与判断能力，而不是取消人的责任。}

\subsection{从 Med-PaLM 到假设生成：为什么不能只靠一个 prompt}

项目的起点来自一条具体反馈。Natarajan 团队早期用 Med-PaLM 处理医学问答与总结，一位 Stanford 教授 Gary Peltz 提出：既然模型吸收了大量科学和医学文本，能否让它提出罕见病的因果基因或新的生物学假设？当时的 PaLM 仍会产生大量 hallucination，团队也不知道怎样可靠地把开放式创意变成可检验研究。

\lecturefigure{slide-003.jpg}{从医学问答模型到 AI co-scientist 的能力路径}{Stanford CS25 V6 Lecture 07 官方录像 00:03:04}

\begin{knowledgebox}{读图：三步路径改变了问题类型}
左侧 Med-PaLM 主要回答定义较清楚的问题；中间阶段开始测试模型能否生成医学假设；右侧才是多智能体 AI co-scientist。关键变化不是模型名称，而是任务从“给定输入、输出答案”变成“在开放空间中提出候选、找反例、比较路线、调用工具，并允许专家持续干预”。问题越开放，单次回答越不足以承载可靠性。
\end{knowledgebox}

\teachervoice{00:02:20--00:05:05，老师回顾当时团队的第一反应是“还没准备好”。PaLM 时代的 hallucination 说明，开放式 hypothesis generation 不能只靠一次采样；需要把批判、比较、验证和人类反馈组织成显式流程。}

\begin{warningbox}{创意与可靠性不是同一个轴}
提高 temperature、一次采样更多答案，可能增加多样性，却不会自动提高可证伪性、事实正确率或实验价值。科研 Agent 必须同时管理探索宽度与证据强度：既不能只重复最安全的已知解释，也不能把陌生、流畅的文本误当成真正新颖的机制。
\end{warningbox}

\subsection{本讲的真实范围：录像没有进入 AMIE 主讲段}

官方视频描述同时宣布 AI co-scientist 与 AI co-physician 项目 AMIE，但实际 1:06:32 录像在 co-scientist 案例结束后直接进入问答。讲义因此只覆盖录像中真正讲授的内容；结尾出现的“让医疗专业能力普遍可及”过渡卡只能证明原计划还有医学协作系统主题，不能被扩写成一堂不存在的 AMIE 课程。

\begin{warningbox}{Source boundary}
课程简介、视频描述和 slide 过渡页可以帮助理解原计划，却不能替代实际讲授证据。本讲不补造 AMIE 架构、评测或病例，也不把后续公开材料倒灌成 2026-05-14 课堂内容。
\end{warningbox}

\subsection{本章小结}

本章把目标从“生成更多科学文本”改写成“提高单位时间内被验证的新知识”。AI co-scientist 的核心对象是研究闭环：问题、假设、批判、排序、实验和更新。项目从 Med-PaLM 的医学问答出发，但开放式发现暴露了单模型、单 prompt 的可靠性不足，因而需要多智能体流程、工具反馈和 scientist-in-the-loop 控制。

